\subsection{对两测度乘积的运算}

命题 10. --- 设 \(g\) （相应地，\(g'\) ）是复函数（或取值于 \(\overline{\mathbf{R}}\) 中的函数），定义在 \(T\) （相应地，\(T'\) ）上。

\begin{enumerate}
\def\labelenumi{\alph{enumi})}
\tightlist
\item
  若 \(g\) （相应地，\(g'\) ）关于 \(\mu\) （相应地，\(\mu'\) ）局部可积，则函数 \(g \otimes g': (t, t') \mapsto g(t)g'(t')\) 关于 \(\nu = \mu \otimes \mu'\) 局部可积，且
\end{enumerate}

\[
(g \cdot \mu) \otimes (g ^ {\prime} \cdot \mu^ {\prime}) = (g \otimes g ^ {\prime}) \cdot (\mu \otimes \mu^ {\prime}).\tag{15}
\]

\begin{enumerate}
\def\labelenumi{\alph{enumi})}
\setcounter{enumi}{1}
\item
  反之，若 \(g \otimes g'\) 是局部 \(\nu\) -可积的，且 \(g'\) 不是局部 \(\mu'\) -可忽略的，则 \(g\) 是局部 \(\mu\) -可积的。
\item
  设 K 和 \(K'\) 分别是 T 和 \(T'\) 的两个紧子集；命题 8 的推论 2 表明函数 \((t, t') \mapsto g(t)g'(t')\varphi_{K \times K'}(t, t')\) （它等于 \((g\varphi_{K}) \otimes (g'\varphi_{K'})\) ）是 \(\nu\) -可积的。因此，\(g \otimes g'\) 是局部 \(\nu\) -可积的。于是立即验证 (15) 的第二端满足乘积测度的特征性质（Ch. III, §4, No.~1, 定理 1）。
\item
  现在设 \(g \otimes g'\) 是局部 \(\nu\) -可积的，且 \(g'\) 不是局部 \(\mu'\) -可忽略的。设 \(\mu_1\) 是满足 \(\mu_1 \leqslant \mu\) 的具紧支集正测度；由于 \(g \otimes g'\) 是 \((\mu_1 \otimes \mu')\) -可测的，故 \(t \mapsto g(t)g'(t')\) 是 \(\mu_1\) -可测的（除去一个局部 \(\mu'\) -可忽略的 \(t'\) 值集外）（命题 2）。
\end{enumerate}

由于 \(g'\) 局部 \(\mu'\) -几乎处处不为零，由此推出 g 是 \(\mu_{1}\) -可测的，再把 \(\mu\) 分解成具紧支集测度族之和，即得 g 是 \(\mu\) -可测的（ \(\S2\) , No.~3, 命题 4 及 \(\S2\) , No.~2, 命题 2）。证明了这一点之后，就可化归到 g 和 \(g'\) \(\geqslant0\) 的情形，如有必要则把 g 和 \(g'\) 换成它们的绝对值。设 K 是 T 的任意紧子集，并设 \(K'\) 是 \(T'\) 的满足 \(\int g'\varphi_{K'}d\mu'\neq0\) 的紧子集。由命题 8，

\[
\left(\int^ {\bullet} g \varphi_ {K} d \mu\right) \left(\int^ {\bullet} g ^ {\prime} \varphi_ {K ^ {\prime}} d \mu^ {\prime}\right) = \iint^ {\bullet} (g \otimes g ^ {\prime}) \varphi_ {K \times K ^ {\prime}} d \mu d \mu^ {\prime} <   + \infty .
\]

于是第一端的第一个因子是有限的，这就完成了证明。

由于 Ch. III, §4, No.~2 的命题 3，本命题推广到复测度。

命题 11. --- 设 \(\pi\) （相应地，\(\pi'\) ）是 T（相应地，T′）到一个局部紧空间 T \(_1\) （相应地，T \(_1'\) ）中的映射。

\begin{enumerate}
\def\labelenumi{\alph{enumi})}
\item
  若 \(\pi\) （相应地，\(\pi'\) ）是 \(\mu\) -逆紧的（相应地，\(\mu'\) -逆紧的），则映射 \(\pi \times \pi'\) 是 \((\mu \otimes \mu')\) -逆紧的，且 \((\pi \times \pi')(\mu \otimes \mu') = \pi(\mu) \otimes \pi'(\mu')\) 。
\item
  反之，若 \(\pi \times \pi'\) 是 \((\mu \otimes \mu')\) -逆紧的且 \(\mu' \neq 0\) ，则 \(\pi\) 是 \(\mu\) -逆紧的。
\item
  因为 \(\pi \times \pi'\) 依 No.~2 命题 3 的推论 1 是 \((\mu \times \mu')\) -可测的。另一方面，若 K（相应地，\(K'\) ）是 \(T_1\) （相应地，\(T_1'\) ）的紧子集，则 \(\overline{\pi}^1(K)\) 和 \(\overline{\pi}'(K')\) 分别关于 \(\mu\) 和 \(\mu'\) 本质可积，因此 \(\overline{\pi}^1(K) \times \overline{\pi}'(K')\) 关于 \(\mu \otimes \mu'\) 本质可积（命题 8 的推论 2）。这就证明了 \(\pi \times \pi'\) 是 \((\mu \times \mu')\) -逆紧的。现在置 \(\mu_1 = \pi(\mu)\) 、\(\mu_1' = \pi'(\mu')\) 、\(\nu_1 = (\pi \times \pi')(\mu \otimes \mu')\) ；对 \(f \in \mathcal{K}(T_1)\) 和 \(f' \in \mathcal{K}(T_1')\) ，有
\end{enumerate}

\[
\begin{array}{r l} \iint f (\pi (t)) f ^ {\prime} (\pi^ {\prime} (t ^ {\prime})) d \mu (t) d \mu^ {\prime} (t ^ {\prime}) \\ & = \left(\int f (\pi (t)) d \mu (t)\right) \left(\int f ^ {\prime} (\pi^ {\prime} (t ^ {\prime})) d \mu^ {\prime} (t ^ {\prime})\right) \end{array}
\]

（命题 8 的推论 2），这就证明了 \(\nu_{1} = \mu_{1}\otimes \mu_{1}^{\prime}\) （Ch. III, §4, No.~1, 定理 1）。

\begin{enumerate}
\def\labelenumi{\alph{enumi})}
\setcounter{enumi}{1}
\tightlist
\item
  现在设 \(\pi \times \pi'\) 是 \(\mu \otimes \mu'\) -逆紧的且 \(\mu' \neq 0\) 。设 \(\mu_1\) 是一个 \(\leqslant \mu\) 的具紧支集测度。由于函数 \(\pi \times \pi'\) 关于 \(\mu_1 \otimes \mu'\) 可测，故映射 \(t \mapsto (\pi(t), \pi'(t'))\) 是 \(\mu\) -可测的（除去 \(t'\) 组成局部 \(\mu'\) -可忽略集的情形外）（No.~2, 命题 2）。由于 \(\mu' \neq 0\) ，由此推出 \(\pi\) 是 \(\mu_1\) -可测的，最终得出 \(\pi\) 是 \(\mu\) -可测的（§2, No.~3, 命题 4 及 §2, No.~2, 命题 2）。剩下要证明 \(\mu^{\bullet}(f \circ \pi) < +\infty\) 对每个函数 \(f \in \mathcal{K}_{+}(\mathrm{T}_{1})\) 成立。若 \(\mu\) 为零，这一性质是明显的。若 \(\mu\) 不为零，则 \(\mu \otimes \mu'\) 也不为零，于是 \((\pi \times \pi') (\mu \otimes \mu') \neq 0\) （§6, No.~2, 命题 2 的推论 1）。由 Ch. III, §4, No.~1 的引理 1，存在两个函数 \(g \in \mathcal{K}_{+}(\mathrm{T}_{1})\) 、\(g' \in \mathcal{K}_{+}(\mathrm{T}_{1}')\) ，使得
\end{enumerate}

\[
\langle (\pi \times \pi^ {\prime}) (\mu \otimes \mu^ {\prime}), g \otimes g ^ {\prime} \rangle \neq 0.
\]

依像测度的定义，这个表达式等于 \(\langle\mu\otimes\mu',(g\circ\pi)\otimes(g'\circ\pi')\rangle\) ，故命题 8 蕴涵 \(\mu^{\bullet}(g'\circ\pi')\neq0\) 。于是由命题 8 以及 §6, No.~2 的命题 2，有

\[
\begin{array}{r l} \left(\int^ {\bullet} (f \circ \pi) d \mu\right) & \left(\int^ {\bullet} (g ^ {\prime} \circ \pi^ {\prime}) d \mu^ {\prime}\right) = \iint^ {\bullet} (f \circ \pi) \otimes (g ^ {\prime} \circ \pi^ {\prime}) d \mu d \mu^ {\prime} \\ & = \iint^ {\bullet} (f \otimes g ^ {\prime}) d ((\pi \times \pi^ {\prime}) (\mu \otimes \mu^ {\prime})) <   + \infty . \end{array}
\]

于是第一端中的第一个积分是有限的，这就完成了证明。

这一结果立即推广到两个复测度的乘积（把该陈述应用于它们的绝对值）。对下列命题同样成立。

命题 12. --- 设 X（相应地，X′）是 T（相应地，T′）的一个局部紧子空间。则诱导测度 \((\mu \otimes \mu')_{X \times X'}\) （在局部紧子空间 \(X \times X'\) ——\(T \times T'\) 的子空间——上）等于乘积 \(\mu_{X} \otimes \mu'_{X'}\) ，即 X 和 \(X'\) 上分别由 \(\mu\) 和 \(\mu'\) 诱导的测度之积。

因为若 \(f \in \mathcal{K}(\mathrm{X})\) 且 \(f' \in \mathcal{K}(\mathrm{X}')\) ，则

\[
\iint_ {\mathrm{X} \times \mathrm{X} ^ {\prime}} f (t) f ^ {\prime} (t ^ {\prime}) d \mu (t) d \mu^ {\prime} (t ^ {\prime}) = \left(\int_ {\mathrm{X}} f (t) d \mu (t)\right) \left(\int_ {\mathrm{X} ^ {\prime}} f ^ {\prime} (t ^ {\prime}) d \mu^ {\prime} (t ^ {\prime})\right)
\]

这依据命题 8 的推论 2，它依诱导测度的定义（Ch. IV, §5, No.~7）证明了

\[
(\mu \otimes \mu^ {\prime}) _ {\mathbf {X} \times \mathbf {X} ^ {\prime}} = \mu_ {\mathbf {X}} \otimes \mu_ {\mathbf {X} ^ {\prime}} ^ {\prime}
\]

（Ch. III, §4, No.~1, 定理 1）。

